
\documentclass{article}
\usepackage{colortbl}
\usepackage{makecell}
\usepackage{multirow}
\usepackage{supertabular}

\begin{document}

\newcounter{utterance}

\centering \large Interaction Transcript for game `taboo', experiment `medium\_en', episode 13 with reka{-}flash{-}3.1{-}t0.0.
\vspace{24pt}

{ \footnotesize  \setcounter{utterance}{1}
\setlength{\tabcolsep}{0pt}
\begin{supertabular}{c@{$\;$}|p{.15\linewidth}@{}p{.15\linewidth}p{.15\linewidth}p{.15\linewidth}p{.15\linewidth}p{.15\linewidth}}
    \# & $\;$A & \multicolumn{4}{c}{Game Master} & $\;\:$B\\
    \hline

    \theutterance \stepcounter{utterance}  
    & & \multicolumn{4}{p{0.6\linewidth}}{
        \cellcolor[rgb]{0.9,0.9,0.9}{
            \makecell[{{p{\linewidth}}}]{
                \texttt{\tiny{[P1$\langle$GM]}}
                \texttt{You are playing a collaborative word guessing game in which you have to describe a target word for another player to guess.} \\
\\ 
\texttt{Rules:} \\
\texttt{(a) You have to reply in the form: CLUE: <some text>. Guesses from the other player will start with GUESS.} \\
\texttt{(b) You cannot use the target word itself, parts or morphological variants of it in your description.} \\
\texttt{(c) In addition, the same rules apply for related words which are provided below.} \\
\\ 
\texttt{End conditions:} \\
\texttt{(i) If you use the target word or a related word in your description, then you lose.} \\
\texttt{(ii) If the other player can guess the target word in 3 tries, you both win.} \\
\\ 
\texttt{Let us start.} \\
\\ 
\texttt{This is the target word that you need to describe and that the other player needs to guess:} \\
\\ 
\texttt{normally} \\
\\ 
\texttt{Related words are:} \\
\\ 
\texttt{{-} unremarkably} \\
\texttt{{-} regularly} \\
\texttt{{-} usually} \\
\\ 
\texttt{Important: You are under time pressure, give short descriptions that are to the point!} \\
            }
        }
    }
    & \\ \\

    \theutterance \stepcounter{utterance}  
    & \multicolumn{4}{p{0.6\linewidth}}{
        \cellcolor[rgb]{0.9,0.9,0.9}{
            \makecell[{{p{\linewidth}}}]{
                \texttt{\tiny{[P1$\rangle$GM]}}
                \texttt{<reasoning>} \\
\texttt{I need to come up with a clue for the word "normally" without using the word itself, any parts of it, or the related words like unremarkably, regularly, usually. Let me think.} \\
\\ 
\texttt{First, "normally" means in a usual or standard way. So maybe something about routine or the usual state. But I can't use "usually" or "regularly". How else can I describe it? Maybe using opposites? Like the opposite of "abnormally" but that's a morphological variant. Wait, the rules say I can't use morphological variants. So "abnormally" is out.} \\
\\ 
\texttt{What about synonyms? Words like "ordinarily", "commonly", but those might be related. The related words given are unremarkably, regularly, usually. So I can't use those. Maybe think of actions that happen typically. Like "as expected" but that's a phrase. Or "the standard way". But "standard" might be okay. Let me check the rules again. The clue can't use the target word or related words. So "standard" isn't a related word here. Maybe "as expected" is a good clue. Let me see. If I say "CLUE: As expected", does that work? It implies that something happens in a usual manner. The other player would guess "normally" because it's the usual way. But does "as expected" violate any rules? It doesn't use the target word or related words. Yeah, that seems safe. Alternatively, "CLUE: The usual way". But "usual" is a related word. Wait, the related words are unremarkably, regularly, usually. So "usually" is a related word. So I can't use "usual" either. Hmm. So "the usual way" is out because of "usual". So maybe "standard" is better. "CLUE: Standard behavior". But "standard" might be too vague. Alternatively, "CLUE: Expected outcome". That could work. If the guesser thinks of the expected way things happen, they might go for "normally". Let me check if "expected" is allowed. The rules don't restrict that. So "CLUE: Expected outcome" is a possible clue. Another angle: "CLUE: Routine activity". But "routine" is similar to regularly, which is a related word. Wait, the related words are unremark} \\
            }
        }
    }
    & & \\ \\

    \theutterance \stepcounter{utterance}  
    & & & \multicolumn{2}{p{0.3\linewidth}}{
        \cellcolor[rgb]{0.9,0.9,0.9}{
            \makecell[{{p{\linewidth}}}]{
                \texttt{\tiny{[GM$|$GM]}}
                \texttt{abort game} \\
            }
        }
    }
    & & \\ \\

\end{supertabular}
}

\end{document}
